\par\clearpage
\normalmarginpar
\scanpage{5}
\mspage{\handwritten{\tikz[baseline=(f.base)]{\node[inner sep=1pt] (f) {11bis};\draw (f.north west) .. controls ([yshift=-5pt]f.south west) and ([yshift=-5pt]f.south) .. ([yshift=-5pt]f.south east);}}}
\noindent\handwritten{\tikz[overlay]{\draw (-1.3,.30)--(-1.24,.41);\draw (-1.14,.41)--(-.99,.29);}}
\handwritten{\tikz[overlay]{\draw (.72\linewidth,.35) .. controls (.73\linewidth,.23) and (.732\linewidth,.23) .. (.733\linewidth,.35);}}
% The three large cancellation strokes cross the two numbered sections.
\figure{\begin{tikzpicture}
\node[anchor=north west,inner sep=0pt] (body) at (0,0) {\begin{minipage}{.93\linewidth}
\small
\noindent
\deleted{1.3.{\deleted{\unsure{2|3}}}{\unsure{3|2}}.}\quad
\deleted{\inserted{The terminology of 1.3.1 {\deleted{and 1.3.2}} extends at once to the case where}}
\deleted{\deleted{Let us return to the case {\unsure{where|or}} one is given two Groups $P$ and $Q$ over $G$}}
\deleted{one has a {\uline{bitorsor}} {\deleted{\unsure{osus|under}}} {\inserted{$E$}} under{\deleted{$(G$}} $(G_{P},G_{P})$, {\deleted{resp{\unsure{.|,}} under $({G\unsure{,}}_{P\mathrm{x}Q},G_{P\mathrm{x}Q})$,}} it being understood that a rigidification of the bitorsor $E$ must be an isomorphism of $e_{P}(E)$ with the trivial b{\deleted{o}}it{\deleted{r}}o{\deleted{o}}rs{\deleted{\unsure{e|u}}}or $\,{}_{s}G_{d}$, {\deleted{thus}} {\inserted{i.e.}}{\handwritten{\tikz[baseline]{\draw (0,-.04) .. controls (.035,.14) and (.16,.06) .. (.19,.25);}}} is defined by a section $e_{E}$ of $e_{P}(E)$ which is {\uline{central}}, i.e. which satisfies the condition $g.e_{E}=e_{E}.g\quad(g\mathrel{\handwritten{\unsure{\in|\epsilon}}}G(S)$ . One similarly extends the terminology of birigidifications to the case of bitorsors.}

\noindent\deleted{1.3.4.}\quad
\deleted{Let us return to the case where one is given two Groups $P$ and {\deleted{$Q$}} {\handwritten{\inserted{$G$}}}, and let us choose the unit section $e_{P}$ to define rigidifications of bitorsors over $P$, and birigidifications of bitorsors over $P\mathrm{x}P$. One has natural functors}

\begingroup
\csname tagsleft@true\endcsname
\diagram{\[
\begin{tikzcd}[ampersand replacement=\&,column sep=1.7em]
\deleted{\underline{\mathrm{Ext}}(P;\unsure{G'|G})} \arrow[r,"{\handwritten{i}}"]
\& \deleted{\underline{\mathrm{Bitorsrig}}(P,G)} \arrow[r,"{\handwritten{\rlap{\raisebox{.35ex}{\hspace{-.12em}\(\color{gray!55}\scriptstyle\delta\)}}\delta}}"]
\& \deleted{\underline{\mathrm{Bitorsbirig}}(P,P;G){,}}
\end{tikzcd}
\tag*{\deleted{(1.3.4.1)}}
\]}
\endgroup

\noindent\deleted{where the three expressions written denote respectively the category of {\handwritten{\tikz[baseline,overlay]{\node[anchor=south west,inner sep=0pt] at (-.15,.31) {\inserted{Group}};\draw (-.08,-.12) .. controls (-.08,.24) and (.06,.26) .. (.35,.26) .. controls (.75,.26) and (1.15,.24) .. (1.35,.32) .. controls (1.40,.35) and (1.40,.43) .. (1.32,.42);}}}extensions of $P$ by $G$, the category of {\deleted{b}}{\handwritten{B}}itorsors under $(G_{P},G_{P})$ rigidified at{\textasciigrave} $e_{P}$, and the category of {\deleted{t}}b{\deleted{o}}it{\deleted{r}}orsors under$(G_{P\mathrm{x}P},G_{P\mathrm{x}P})$ birigidified at $(e_{P},e_{P})$. One {\deleted{takes}} defines the functor $i$ by associating to the extension $E$ of $P$ by $G$ the corresponding {\deleted{torso}} bitorsor over $P$ (1.1.6), rigidified by the unit sectio{ń} of $E$. O{\deleted{\unsure{n}}}{\handwritten{n}}e {\unsure{defines|defines}} the functor $\handwritten{\delta}$ by the formu{\deleted{o}}la}
\end{minipage}};
\begin{scope}[x={(body.north east)},y={(body.south west)}]
\draw (.44,.11) .. controls (.35,.27) and (.00,.66) .. (-.047,.725)
 .. controls (-.053,.736) and (-.054,.746) .. (-.043,.75);
\draw (.82,.078) .. controls (.73,.30) and (.50,.69) .. (.31,1.035);
\draw (.91,.505) .. controls (.915,.62) and (.78,.86) .. (.655,1.03);
\end{scope}
\end{tikzpicture}}
% archive sheet number: 4
